% Part: proof-theory
% Chapter: cut-elimination
% Section: midsequent

\documentclass[../../../include/open-logic-section]{subfiles}

\begin{document}

\olfileid[hi]{pt}{cut}{mid}

\olsection{मध्य-अनुक्रम प्रमेय और हेरब्रांड प्रमेय}

कट-विलोपन प्रमेय का एक महत्त्वपूर्ण परिणाम मध्य-अनुक्रम प्रमेय है। यह कहता है
कि यदि~$\Gamma \Sequent \Delta$ का ऐसा !!a{proof}~$\pi$ है जिसमें प्रत्येक
!!{formula} पूर्वबंध रूप में है, तो ऐसा !!a{proof}~$\pi'$ है जिसमें कोई अनुक्रम
$S \ident \Pi \Sequent \Lambda$ (जिसे \emph{मध्य-अनुक्रम} कहते हैं) है और
~$\pi'$ में $S$ के ऊपर के सभी अनुमान प्रतिज्ञप्तिपरक तथा~$S$ के नीचे के सभी
अनुमान परिमाणक अनुमान हैं। (!!^a{formula}~$!A$ \emph{पूर्वबंध} (या
\emph{पूर्वबंध रूप में}) है यदि उसका रूप
\[\mathsf{Q}_1x_1\dots\mathsf{Q}_nx_n\,!B(x_1,\dots,x_n)\] है, जहाँ प्रत्येक
$\mathsf{Q}_i$ या तो $\lforall$ या~$\lexists$ है।) मध्य-अनुक्रम प्रमेय का एक
महत्त्वपूर्ण परिणाम हेरब्रांड प्रमेय है। उसकी सामान्य स्थिति का वर्णन कुछ
कठिन है, पर बहुत व्यापक रूप में वह यह कहता है: प्रथम-कोटि तर्क के प्रत्येक
!!{provable} !!{sentence}~$!A$ के लिए कोई प्रतिज्ञप्तिपरक सर्वतासत्य~$!A'$ है
जिससे $!A$ को !!{prove} किया जा सकता है। यहाँ उस स्थिति का रूप है जहाँ
!!{sentence}~$\lexists[x][!B(x)]$ है और $!B$ परिमाणक-मुक्त है:

\begin{thm}[हेरब्रांड प्रमेय]
मान लें $!B(x)$ परिमाणक-मुक्त है और $\Proves \lexists[x][!B(x)]$। तब ऐसे
पद~$t_1$, \dots, $t_n$ हैं कि $\Proves !B(t_1) \lor \dots \lor !B(t_n)$।
\end{thm}

!!{formula}~$!B(t_1) \lor \dots \lor !B(t_n)$ को~$\lexists[x][!B(x)]$ का
\emph{हेरब्रांड विकल्प} कहते हैं। चूँकि $!B$ परिमाणक-मुक्त है, यह विकल्प भी
परिमाणक-मुक्त है। अतः यदि वह !!{provable} है, तो वह किसी कट-मुक्त प्रमाण से,
और फलतः केवल प्रतिज्ञप्तिपरक अनुमानों का उपयोग करने वाले !!a{proof} से,
!!{provable} है। इसलिए वह सर्वतासत्य है।

हेरब्रांड प्रमेय का कठिन भाग यह दिखाना है कि प्रत्येक !!{provable}
!!{formula}~$\lexists[x][!B(x)]$ के लिए कोई हेरब्रांड विकल्प विद्यमान है।
इसका विलोम—यदि $\lexists[x][!B(x)]$ का कोई हेरब्रांड विकल्प है तो वह
सिद्ध-योग्य है—सरल है। वास्तव में हेरब्रांड विकल्प, $\lexists[x][!B(x)]$ को
निगमित करता है। उदाहरणतः $n=2$ वाली स्थिति का~$\Log{G3c}$ में !!a{proof} यह
है:
\[
\Axiom$!B(t_1) \fCenter \lexists[x][!B(x)], !B(t_1), !B(t_2)$
\Axiom$!B(t_2) \fCenter \lexists[x][!B(x)], !B(t_1), !B(t_2)$
\RightLabel{\LeftR{\lor}}
\BinaryInf$!B(t_1) \lor !B(t_2) \fCenter \lexists[x][!B(x)], !B(t_1), !B(t_2)$
\RightLabel{\RightR{\lexists}}
\UnaryInf$!B(t_1) \lor !B(t_2) \fCenter \lexists[x][!B(x)], !B(t_2)$
\RightLabel{\RightR{\lexists}}
\UnaryInf$!B(t_1) \lor !B(t_2) \fCenter \lexists[x][!B(x)]$
\DisplayProof
\]

~$\Sequent \lexists[x][!B(x)]$ के !!a{proof}~$\pi$ से हेरब्रांड विकल्प निकालने
के लिए उस स्थिति पर विचार करें जहाँ कट-मुक्त !!{proof}~$\pi$ के सभी
$\RightR{\lexists}$ अनुमान अंत में हैं:
\[
\AxiomC{}
\RightLabel{$\pi_1$}
\Deduce$ \fCenter \lexists[x][!B(x)], !B(t_1), \dots, !B(t_n)$
\RightLabel{\RightR{\lexists}}
\UnaryInf$ \fCenter \lexists[x][!B(x)], !B(t_2), \dots, !B(t_n)$
\RightLabel{$(n-1) \times \RightR{\lexists}$}
\Deduce$ \fCenter \lexists[x][!B(x)]$
\DisplayProof
\]
यदि ये~$\pi$ के वे सभी $\RightR{\lexists}$ अनुमान हैं जिनमें
$\lexists[x][!B(x)]$ सक्रिय है, तो~$\pi_1'$ में कोई अन्य परिमाणक अनुमान नहीं
है। कारण यह है कि $!B(x)$ में कोई अन्य परिमाणक नहीं और अंतिम अनुक्रम की बाईं
ओर कोई !!{formula} नहीं है। दूसरे शब्दों में, $ \fCenter
\lexists[x][!B(x)], !B(t_1), \dots, !B(t_n)$,~$\pi$ का मध्य-अनुक्रम है। इसके
अतिरिक्त, मध्य-अनुक्रम के ऊपर कोई परिमाणक अनुमान न होने से उसके ऊपर के सभी
अनुक्रमों में $\lexists[x][!B(x)]$ है, और ये न सक्रिय हैं न प्रमुख। अतः इन्हें
!!{proof}~$\pi_1$ से हटाया जा सकता है; इससे~$\fCenter !B(t_1), \dots, !B(t_n)$
का !!a{proof}, और $\RightR{\lor}$ के $n-1$ अनुप्रयोगों से हेरब्रांड विकल्प
~$\Sequent !B(t_1) \lor \dots \lor !B(t_n)$ का !!a{proof}, मिलता है।

समस्या यह है कि हम यह नहीं मान सकते कि $\Sequent \lexists[x][!B(x)]$ के
कट-मुक्त !!{proof} में भी सभी \RightR{\lexists} अनुमान अंत में एक ही खंड में
होंगे। इन \RightR{\lexists} अनुमानों के बीच ऐसे अनुमान हो सकते हैं जिनमें कोई
$!B(t_i)$ प्रमुख हो। इसलिए मध्य-अनुक्रम प्रमेय की आवश्यकता पड़ती है।

\begin{thm}[मध्य-अनुक्रम प्रमेय]\ollabel{thm:midsequent} यदि $\Gamma \Sequent
  \Delta$ का कोई \Log{G3c}-!!{proof}~$\pi$ है जिसमें $\Gamma$ और $\Delta$ के
  सभी !!{formula}s पूर्वबंध हैं, तो ऐसा !!a{proof}~$\pi'$ है जिसमें कोई अनुक्रम
  $\Pi \Sequent \Lambda$ है और उसके नीचे के सभी अनुमान परिमाणक अनुमान तथा
  उसके ऊपर के सभी अनुमान प्रतिज्ञप्तिपरक अनुमान हैं।
\end{thm}

\begin{proof}
  ~ $\pi$ में किसी परिमाणक अनुमान की \emph{कोटि} उसके नीचे के
  प्रतिज्ञप्तिपरक अनुमानों की संख्या, और $\pi$ की कोटि~$o(\pi)$ उसके सभी
  परिमाणक अनुमानों की कोटियों का योग परिभाषित करें। यदि $o(\pi) = 0$, तो किसी
  परिमाणक अनुमान के नीचे कोई प्रतिज्ञप्तिपरक अनुमान नहीं है और काम पूरा हुआ:
  सभी परिमाणक अनुमानों का एक आधार-वाक्य होने से केवल एक सर्वोच्च परिमाणक
  अनुमान है। उसका आधार-वाक्य मध्य-अनुक्रम है।

  यदि $o(\pi) > 0$, तो कोटि $>0$ वाला कोई सबसे निचला परिमाणक अनुमान चुनें।
  उसकी कोटि $>0$ होने से उसके नीचे कोई प्रतिज्ञप्तिपरक अनुमान अवश्य है। उसके
  सबसे निचला होने से उसका निष्कर्ष किसी परिमाणक अनुमान का आधार-वाक्य नहीं हो
  सकता: अन्यथा उस दूसरे, अधिक नीचे वाले परिमाणक अनुमान के नीचे भी कोई
  प्रतिज्ञप्तिपरक अनुमान होता। हमारे पास कौन-सा परिमाणक अनुमान है और उसके नीचे
  कौन-सा प्रतिज्ञप्तिपरक अनुमान है, इसके अनुसार हम (बहुत-सी!) स्थितियों में
  भेद करते हैं। अंतिम अनुक्रम में केवल पूर्वबंध सूत्र होने के कारण परिमाणक
  अनुमान का प्रमुख सूत्र अगले प्रतिज्ञप्तिपरक अनुमान में सक्रिय नहीं हो सकता।
  अन्यथा उस अनुमान के प्रमुख सूत्र में किसी प्रतिज्ञप्तिपरक संयोजक के प्रभाव
  क्षेत्र में परिमाणक होता। उप-!!{formula} गुण से उसे अंतिम अनुक्रम के किसी
  !!{formula} का उप-!!{formula} होना पड़ता। अतः परिमाणक अनुमान का प्रमुख
  !!{formula}, नीचे वाले प्रतिज्ञप्तिपरक अनुमान के प्रसंग का !!a{element} है।

  ये दो उदाहरण हैं:

पहले मान लें कोई $\RightR{\lforall}$, $\LeftR{\lif}$ के दाएँ आधार-वाक्य पर
समाप्त होता है। तब~$\pi$ का संबंधित उप-!!{proof}, उप-!!{proof}~$\pi_1$ पर
समाप्त होता है।
\[
\AxiomC{}
\RightLabel{$\pi_2$}
\Deduce$\Gamma' \fCenter \Delta', \lforall[x][!A(x)], !B$
\AxiomC{}
\RightLabel{$\pi_3$}
\Deduce$!C, \Gamma' \fCenter \Delta', !A(c)$
\RightLabel{\RightR{\lforall}}
\UnaryInf$!C, \Gamma' \fCenter \Delta', \lforall[x][!A(x)]$
\RightLabel{\LeftR{\lif}}
\BinaryInf$!B \lif !C, \Gamma' \fCenter \Delta', \lforall[x][!A(x)]$
\DisplayProof
\]
हम $\pi$ को नियमित मान सकते हैं, इसलिए आइगेनचर~$c$,~$\pi_3$ के बाहर नहीं आता;
विशेषतः वह $!B \lif !C$ या~$\pi_2$ में नहीं आता। अब प्रमाण~$\pi_1'$ पर विचार
करें:
\[
\AxiomC{}
\RightLabel{$\pi_2'$}
\Deduce$\Gamma' \fCenter \Delta', \lforall[x][!A(x)], !A(c), !B$
\AxiomC{}
\RightLabel{$\pi_3$}
\Deduce$!C, \Gamma' \fCenter \Delta', \lforall[x][!A(x)], !A(c)$
\RightLabel{\LeftR{\lif}}
\BinaryInf$!B \lif !C, \Gamma' \fCenter \Delta', \lforall[x][!A(x)], !A(c)$
\RightLabel{\RightR{\lforall}}
\UnaryInf$!B \lif !C, \Gamma' \fCenter \Delta', \lforall[x][!A(x)], \lforall[x][!A(x)]$
\DisplayProof
\]
जहाँ $\pi_2'$ को~$\pi_2$ से दाईं ओर~$!A(c)$ द्वारा दुर्बलीकरण करके प्राप्त
किया जाता है। (इससे कोई अनुमान नहीं जुड़ता, विशेषतः ऐसा कोई परिमाणक अनुमान
नहीं जो कोटि को प्रभावित कर सके। इससे~$\pi_2$ की किसी आइगेनचर-शर्त का उल्लंघन
भी नहीं होता, क्योंकि~$!A(c)$ का कोई अचर~$\pi_2$ में आइगेनचर के रूप में
प्रयुक्त नहीं है।) \RightR{\lforall} की आइगेनचर-शर्त अब भी पूरी है, क्योंकि
$c$,~$!B \lif !C$ में नहीं आता।

संकुचन की ग्राह्यता का आह्वान करके हम $!B \lif !C, \Gamma' \fCenter \Delta',
\lforall[x][!A(x)]$ का !!a{proof}~$\pi_1''$ प्राप्त करते हैं।
$\lforall[x][!A(x)]$ के लिए $\RightR{\Contraction}$ की ग्राह्यता के प्रमाण और
उसमें प्रयुक्त $\RightR{\lforall}$ की उत्क्रमणीयता के प्रमाण ने किसी अनुमान की
कोटि नहीं बदली: उसने अधिक-से-अधिक अनुमान हटाए। अतः~$\pi_1''$ में परिमाणक
अनुमानों की संख्या~$\pi_1$ से अधिक नहीं है, और~$\pi_1''$ के प्रत्येक परिमाणक
अनुमान के नीचे प्रतिज्ञप्तिपरक अनुमानों की संख्या~$\pi_1$ के संगत अनुमान जैसी
है—केवल अंतिम \RightR{\lforall} को छोड़कर:~$\pi$ में उसके नीचे एक
\LeftR{\lif} था, जिसे हमने~$\pi_1''$ में उसके ऊपर पहुँचा दिया है।~$\pi$ में
$\pi_1$ को $\pi_1''$ से बदलें। परिणामी !!{proof} की कोटि कम है, क्योंकि
कम-से-कम $\RightR{\lforall}$ अनुमान की कोटि (कम-से-कम)~$1$ घटी है। अब आगमन
परिकल्पना लगाएँ।

जहाँ परिमाणक अनुमान \RightR{\lexists} है वे स्थितियाँ और भी सरल हैं, क्योंकि
न संकुचन चाहिए, न किसी आइगेनचर-शर्त की चिंता। उदाहरणतः मान लें
$\RightR{\lexists}$ के बाद $\RightR{\land}$ आता है (इस बार बाएँ आधार-वाक्य के
रूप में)। संबंधित उप-!!{proof} $\pi_1$ है:
\[
\AxiomC{}
\RightLabel{$\pi_2$}
\Deduce$\Gamma' \fCenter \Delta', \lexists[x][!A(x)], !A(t), !B$
\RightLabel{\RightR{\lexists}}
\UnaryInf$!\Gamma' \fCenter \Delta', \lexists[x][!A(x)], !B$
\AxiomC{}
\RightLabel{$\pi_3$}
\Deduce$\Gamma' \fCenter \Delta', \lexists[x][!A(x)], !C$
\RightLabel{\RightR{\land}}
\BinaryInf$\Gamma' \fCenter \Delta', \lexists[x][!A(x)], !B \land !C$
\DisplayProof
\]
~$\pi$ में $\pi_1$ को $\pi_1'$ से बदलें:
\[
\AxiomC{}
\RightLabel{$\pi_2$}
\Deduce$\Gamma' \fCenter \Delta', \lexists[x][!A(x)], !A(t), !B$
\AxiomC{}
\RightLabel{$\pi_3'$}
\Deduce$\Gamma' \fCenter \Delta', \lexists[x][!A(x)], !A(t), !C$
\RightLabel{\RightR{\land}}
\BinaryInf$\Gamma' \fCenter \Delta', \lexists[x][!A(x)], !A(t), !B \land !C$
\RightLabel{\RightR{\lexists}}
\UnaryInf$\Gamma' \fCenter \Delta', \lexists[x][!A(x)], !B \land !C$
\DisplayProof
\]
जहाँ $\pi_3'$ को $\pi_3$ से दाईं ओर~$!A(t)$ द्वारा दुर्बलीकरण करके प्राप्त
किया जाता है। इस दुर्बलीकरण में केवल~$\pi_3$ के प्रत्येक अनुक्रम की दाईं ओर
$!A(t)$ जोड़ना है, इसलिए किसी परिमाणक अनुमान की कोटि नहीं बदलती।~$\pi'$ में
परिमाणक अनुमानों की कोटियाँ~$\pi$ जैसी हैं, केवल~\RightR{\land} के साथ
क्रम-विनिमय किए गए \RightR{\lexists} अनुमान की कोटि~$1$ घटी है। आगमन
परिकल्पना लागू होती है।
\end{proof}

\begin{prob}
\olref{thm:midsequent} के प्रमाण की उन स्थितियों का वर्णन करें जहाँ
\LeftR{\lexists},~\RightR{\lif} के आधार-वाक्य पर समाप्त होता है, और जहाँ
\LeftR{\lforall},~\LeftR{\lor} के बाएँ आधार-वाक्य पर समाप्त होता है।
\end{prob}

\begin{proof}[हेरब्रांड प्रमेय का प्रमाण]
मान लें $\pi$, $\Sequent \lexists[x][!A(x)]$ पर समाप्त होता है।
\olref{thm:midsequent} से हम $\pi$ को किसी मध्य-अनुक्रम~$S$ वाले
!!a{proof}~$\pi'$ में बदल सकते हैं, जिसका रूप अवश्य \[\Sequent
\lexists[x][!A(x)], !A(t_1), \dots, !A(t_n)\] है। चूँकि $S$ के ऊपर के सभी
अनुमान प्रतिज्ञप्तिपरक हैं, $\lexists[x][!A(x)]$,~$S$ के ऊपर किसी अनुमान में
प्रमुख नहीं हो सकता। परमाण्विक न होने से वह किसी अभिगृहीत में प्रमुख नहीं हो
सकता। चूँकि $\lexists[x][!A(x)]$, $!A(t_1)$ का उचित उपसूत्र नहीं हो सकता (याद
रखें कि $!A(x)$ परिमाणक-मुक्त है), वह~$S$ के ऊपर सक्रिय भी नहीं हो सकता। अतः
~$S$ पर समाप्त होने वाले उप-!!{proof} से $\lexists[x][!A(x)]$ हटाने पर
$\Sequent !A(t_1), \dots, !A(t_n)$ का सही !!{proof} मिलता है; उससे $n-1$
\RightR{\lor} अनुमानों द्वारा $\Sequent !A(t_1) \lor \dots \lor !A(t_n)$ का
!!a{proof} प्राप्त होता है।
\end{proof}

% \begin{prob}
% Find a Herbrand disjunction of $\lexists[x][\lforall[y][(!A(x) \lif
% !A(y))]]$ by finding a cut-free !!{proof} of $\Sequent
% \lexists[x][\lforall[y][(!A(x) \lif !A(y))]]$ and applying the
% transformation in \olref{thm:midsequent} to find a midsequent (if
% necessary).
% \end{prob}

\end{document}
